\section{The Dying Begun in Baptism}

Paul speaks to the baptized as people whose relation to death has already changed. They have been buried with Christ through Baptism so that they may walk in newness of life. Their union with his death includes the promise of union with his Resurrection (Rom 6:3--11). Through the sacrament we really participate in his saving mystery. We receive a life that begins now and will be fulfilled when he raises us from the dead.

Paul immediately joins this gift to conduct. The baptized are to refuse sin's reign in their mortal bodies and offer themselves to God as people brought from death to life (Rom 6:12--13). The body that will one day die is presently an instrument of obedience. Its hands, voice, eyes, strength, and time can be offered to charity. Preparation for death therefore includes how we inhabit the body now: the uses we make of it, the restraints we accept, the persons we serve, and the worship in which we place ourselves before God.

Baptism establishes our life in Christ before we have done anything to earn it. Many of us were carried to the font before we could form an intention of our own. The Church confessed the faith over us, and the Lord received us through her ministry. At the other end of life we may again need to be carried, spoken for, and surrounded by prayer. Receiving this help can remind us how much of our Christian life has depended on others: those who brought us to Baptism, taught us to pray, corrected us, or helped us return after we had fallen away.

We must freely respond to the grace we have received. We belong to Christ, and our choices should accord with that gift. Dying to self takes place in the ordinary duties through which we learn to obey him. We keep a promise when inconvenience has made it costly. We accept a correction without immediately preparing a defense. We undertake a duty that brings no admiration. We refuse an indulgence because it is injuring our freedom to love. In each case, the sacrifice concerns something we have wanted to keep at the expense of fidelity.

None of this requires contempt for created goods. Food, friendship, marriage, work, and rest are gifts through which we can give thanks to God. Baptismal renunciation opposes their disorder, the possessiveness that makes a gift rival its Giver. We can learn to enjoy a good thing without demanding that it last forever in its present form. We can work carefully without claiming final control of the result. We can love a person deeply without making that person responsible for providing the security that belongs to God alone.

This education in freedom becomes especially visible when a familiar good is withdrawn. A change in health may reveal how much self-respect depended on strength. Retirement may expose a need to be necessary. A child's departure may show where parental love has become a demand to direct every outcome. These revelations can be painful, but they offer places for grace to work. The difficulty is no proof that earlier love or work was worthless. It reveals an attachment within something good that the Lord is teaching us to surrender more freely.

Fasting, almsgiving, obedient service, and patient prayer help us submit to Christ in the use of these goods. Fasting disciplines an appetite; almsgiving makes possessions available to someone in need; service requires time we might have kept for ourselves. Their Christian purpose is fulfilled as they help us love God and our neighbor more freely. When death removes the bodily powers by which we have practiced them, we still rely on the life received in Baptism. The Lord with whom we have died sacramentally has promised to raise us with him.

